\BourbakiExerciseGroup

\begin{enumerate}
\def\labelenumi{\arabic{enumi})}
\tightlist
\item
  \begin{enumerate}
  \def\labelenumii{\alph{enumii})}
  \tightlist
  \item
    Let K be a field, f a separable polynomial in K{[}X{]} and N a splitting field of f.~Show that for the group \(\operatorname{Gal}(N/K)\) to operate transitively on the set of roots of f in N it is necessary and sufficient for f to be irreducible in K{[}X{]}.
  \end{enumerate}
\end{enumerate}

\begin{enumerate}
\def\labelenumi{\alph{enumi})}
\setcounter{enumi}{1}
\tightlist
\item
  Suppose further that \(f\) is irreducible in \(\mathbf{K}[\mathbf{X}]\) ; let \(x\) be a root of \(f\) in \(\mathbf{N}\) . For there to exist no field \(\mathbf{E}\) such that \(\mathbf{K} \subset \mathbf{E} \subset \mathbf{K}(x)\) , distinct from \(\mathbf{K}\) and \(\mathbf{K}(x)\) it is necessary and sufficient that \(\operatorname{Gal}(\mathbf{N} / \mathbf{K})\) , qua transitive permutation group of the set of roots of \(f\) in \(\mathbf{N}\) , should be primitive (I, p.~142, Ex. 13).
\end{enumerate}

\begin{enumerate}
\def\labelenumi{\arabic{enumi})}
\setcounter{enumi}{1}
\tightlist
\item
  Let K be a field, \(f_{0}\) an irreducible and separable polynomial in K{[}X{]}, of degree n; let \(\alpha_{i}\) ( \(1 \leqslant i \leqslant n\) ) be its roots in a splitting field N of \(f_{0}\) . Let F be the field of rational fractions \(\mathrm{N}(\mathbf{X}_{1}, \mathbf{X}_{2}, \ldots, \mathbf{X}_{n})\) , E the subfield \(\mathrm{K}(\mathbf{X}_{1}, \mathbf{X}_{2}, \ldots, \mathbf{X}_{n})\) of F; F is a Galois extension of E and \(\operatorname{Gal}(\mathrm{E}/\mathrm{F})\) is canonically isomorphic to \(\operatorname{Gal}(\mathrm{N}/\mathrm{K})\) (V, p.~71, Th. 5; V, p.~154, Ex. 8). Show that if \(\theta = \alpha_{1}\mathbf{X}_{1} + \alpha_{2}\mathbf{X}_{2} + \cdots + \alpha_{n}\mathbf{X}_{n}\) then \(\mathrm{F} = \mathrm{E}(\theta)\) and the minimal polynomial \(g_{1}\) of \(\theta\) in E{[}T{]} is an irreducible factor of the polynomial
\end{enumerate}

\[
f (T) = \prod_ {\sigma \in \mathfrak {S} _ {n}} \left(T - \alpha_ {1} X _ {\sigma (1)} - \alpha_ {2} X _ {\sigma (2)} - \dots - \alpha_ {n} X _ {\sigma (n)}\right).
\]

Each permutation \(\sigma \in \mathfrak{S}_n\) defines a K-automorphism of E, again denoted by \(\sigma\) , by the condition \(\sigma(X_i) = X_{\sigma(i)}\) for \(1 \leqslant i \leqslant n\) ; show that \(\operatorname{Gal}(N/K)\) is isomorphic to the subgroup of \(\mathfrak{S}_n\) consisting of permutations \(\sigma\) such that \(\sigma(g_1) = g_1\) . Further, if \(f = g_1g_2 \ldots g_r\) is a decomposition of \(f\) into irreducible factors in E{[}T{]}, show that for every index \(j\) such that \(1 \leqslant j \leqslant r\) there exists a permutation \(\sigma_j \in \mathfrak{S}_n\) such that \(\sigma_j(g_1) = g_j\) .

\begin{enumerate}
\def\labelenumi{\arabic{enumi})}
\setcounter{enumi}{2}
\item
  Let N be a Galois extension of a field K of infinite degree over K. Show that the cardinal of the Galois group \(\operatorname{Gal}(N/K)\) is greater than that of the set of subsets of N (use V, p.~61 and 62). Deduce that there exist non-closed subgroups of \(\operatorname{Gal}(N/K)\) .
\item
  Let \(\mathbf{N}\) be a Galois extension of a field \(\mathbf{K}\) , \((\mathrm{E}_i)_{i\in I}\) a family of subextensions of \(\mathbf{N}\) , Galois over \(\mathbf{K}\) , such that: \(1^{\circ}\) for each index \(i\) , if \(\mathbf{F}_i\) is the field generated by the \(\mathbf{E}_j\) for \(j\neq i\) , then \(\mathbf{E}_i\cap \mathbf{F}_i = \mathbf{K}\) ; \(2^{\circ}\) \(\mathbf{N}\) is generated by the union of the \(\mathbf{E}_i\) .
\end{enumerate}

\begin{enumerate}
\def\labelenumi{\alph{enumi})}
\tightlist
\item
  Show that \(\mathbf{N}\) is isomorphic to the tensor product \(\otimes_{i\in I}\mathrm{E}_i\) (III, p.~470) (define an
\end{enumerate}

isomorphism of this tensor product onto N, using Th. 5 of V, p.~71).

\begin{enumerate}
\def\labelenumi{\alph{enumi})}
\setcounter{enumi}{1}
\tightlist
\item
  Show that \(\operatorname{Gal}(\mathbf{N}/\mathbf{K})\) is isomorphic to the product of the topological groups \(\operatorname{Gal}(\operatorname{E}_{i}/\operatorname{K})\) .
\end{enumerate}

\begin{enumerate}
\def\labelenumi{\arabic{enumi})}
\setcounter{enumi}{4}
\item
  Show that every compact totally disconnected group is isomorphic to the group Gal(N/K) of a Galois extension of a field (using Gen.~Top. III, p.~325, Ex. 3, reduce to the case where the group is a product of finite groups and use Ex. 4 as well as V, p.~59, Example 5).
\item
  Let N be a Galois extension of a field K, of infinite degree, and let \(\rho, \sigma\) be two distinct elements of \(\operatorname{Gal}(N/K)\) . If J is the set of \(x \in N\) such that \(\rho(x) \neq \sigma(x)\) , show that \(\mathrm{K}(\mathrm{J})\) is an extension of K of infinite degree. (Observe that a field cannot be the union of two of its subfields unless it is equal to one of them.)
\item
  Let \(\mathbf{N}\) be a Galois extension of a field \(\mathbf{K}\) such that \(\operatorname{Gal}(\mathbf{N} / \mathbf{K})\) is isomorphic to the symmetric group \(\mathfrak{S}_n\) (V, p.~59, Example 5), \(n\) being a non-prime integer; identify \(\operatorname{Gal}(\mathbf{N} / \mathbf{K})\) and \(\mathfrak{S}_n\) . Let \(\Gamma_1\) be the subgroup of \(\mathfrak{S}_n\) of order \((n - 1)!\) leaving the number 1 invariant and let \(\Gamma_2\) be the cyclic subgroup of \(\mathfrak{S}_n\) of order \(n\) generated by the cycle \((1, 2, \ldots, n)\) (I, p.~142, Ex. 12). Let \(E_1, E_2\) be the field of invariants of the subgroups \(\Gamma_1, \Gamma_2\) respectively. Show that \(E_1\) and \(E_2\) are linearly disjoint over \(\mathbf{K}\) , that there exists no field \(F\) other than \(\mathbf{K}\) and \(E_1\) such that \(K \subset F \subset E_1\) but that there exist fields \(F'\) distinct from \(E_2\) and \(N\) such that \(E_2 \subset F' \subset N\) (use Ex. 13 of I, p.~142 and Ex. 14, I, p.~143).
\item
  Let K be a field and N a Galois extension of K of finite degree n.
\end{enumerate}

\begin{enumerate}
\def\labelenumi{\alph{enumi})}
\item
  Let \(p\) be a prime number dividing \(n\) and let \(p^r\) be the highest power of \(p\) dividing \(n\) . Show that there exist subextensions \(L_i\) of \(N\) for \(0 \leqslant i \leqslant r\) such that for \(0 \leqslant i \leqslant r - 1\) , \(L_i\) contains \(L_{i+1}\) , \(L_0 = N\) , \(L_i\) is a Galois extension of degree \(p\) of \(L_{i+1}\) and \([L_r: K]\) is not divisible by \(p\) (cf.~I, p.~77, Th. 1 and I, p.~78, Th. 2).
\item
  Let \(x \in \mathbf{N}\) be such that \(\mathbf{N}\) is generated by \(x\) and its conjugates. Show that there exists a subextension \(\mathbf{E}\) of \(\mathbf{N}\) such that \(x \notin \mathbf{E}\) and \([\mathbf{N} : \mathbf{E}] = p'\) . (Consider the Sylow \(p\) -groups \(\mathrm{H}_j\) of \(\operatorname{Gal}(\mathbf{N} / \mathbf{K})\) and the (normal) subgroup \(\mathbf{H}\) generated by their union; if \(F_j\) (resp. \(\mathbf{F}\) ) is the subfield of invariants of \(\mathrm{H}_j\) (resp. \(\mathbf{H}\) ) show that \(x \notin \mathbf{F}\) and deduce that there is an index \(j\) such that \(x \notin F_j\) ; then we can take \(E = F_{j}\) .) Conclude that there exists a subextension \(L\) of \(\mathbf{N}\) such that \(L(x)\) is a Galois extension of degree \(p\) of \(L\) (use I, p.~77, Prop. 12).
\item
  Conversely, let \(\Omega\) be an algebraic closure of \(N\) and suppose that there exists a subextension \(L\) of \(\Omega\) such that \(L(x)\) is Galois of degree \(p\) over \(L\) . Show that then \(p\) divides \(n\) . (Observe that \(N \cap L(x) = (N \cap L)(x)\) and that \((N \cap L)(x)\) is a Galois extension of degree \(p\) of \(N \cap L\) .)
\end{enumerate}

\begin{enumerate}
\def\labelenumi{\arabic{enumi})}
\setcounter{enumi}{8}
\tightlist
\item
  Let K be a field of characteristic \(p\) , \(\Omega\) an algebraically closed extension of K, \(x\) an element of \(\Omega\) such that \(x \notin K\) and M a maximal element of the set of subextensions of \(\Omega\) not containing \(x\) (V, p.~148, Ex. 9), so that \(\Omega\) is an algebraic extension of M (loc. cit.). a) Show that if M is not a perfect field, \(\Omega\) is a \(p\) -radical extension of M, union of the fields \(M(x^{p^{-n}})\) for \(n \geqslant 0\) (observe that for every element \(y \in \Omega\) \(p\) -radical over M, the field \(M(y)\) contains \(M(x)\) ); deduce that we necessarily have \([M(x):M] = p\) , because every element of \(\Omega\) separable over M is necessarily in M).
\end{enumerate}

\begin{enumerate}
\def\labelenumi{\alph{enumi})}
\setcounter{enumi}{1}
\tightlist
\item
  Show that if M is perfect, \(\mathbf{M}(x)\) is a Galois extension of M, of prime degree q, and for every subextension N of \(\Omega\) which is Galois of finite degree over M, the degree \([N:M]\) is a power of q. (Take for q a prime number dividing \([\mathbf{M}(x):\mathbf{M}]\) ; note that every extension of M distinct from M and contained in \(\Omega\) contains \(\mathbf{M}(x)\) , and use Ex. 8, a).)
\end{enumerate}

\begin{enumerate}
\def\labelenumi{\arabic{enumi})}
\setcounter{enumi}{9}
\tightlist
\item
  Let \(\mathbf{K}\) be a perfect field such that for every integer \(n > 1\) there exists a Galois extension \(\mathbf{N}\) of \(\mathbf{K}\) such that \(\operatorname{Gal}(\mathbf{N} / \mathbf{K})\) is isomorphic to the alternating group \(\mathfrak{A}_n\) ; for example, for every field \(k\) of characteristic 0 the field \(\mathbf{K} = k(\mathbf{X}_n)_{n \in \mathbb{N}}\) of rational fractions in an infinity of indeterminates has this property (V, p.~59, Example 5).
\end{enumerate}

\begin{enumerate}
\def\labelenumi{\alph{enumi})}
\item
  Let A be a finite set of integers \textgreater{} 1 and let \(\Omega\) be an algebraic closure of K. An element \(x \in \Omega\) is said to be A-constructible if there exists an increasing sequence \(K = E_{0} \subset E_{1} \subset \cdots \subset E_{r} \subset \Omega\) of extensions of K such that each of the degrees \([E_{j}: E_{j-1}]\) belongs to A for \(1 \leqslant j \leqslant r\) and \(x \in E_{r}\) . Show that there exists a maximal element L in the set of subextensions \(E \subset \Omega\) such that each \(x \in E\) is A-constructible. For each \(x \notin L\) show that the degree \([L(x): L]\) belongs to N - A.
\item
  Let \(a\) be the greatest element of \(\mathbf{A}\) and let \(n\) be an integer such that \(n \geqslant \max (a + 1,5)\) . Show that there exists an extension \(\mathbf{M}\) of \(\mathbf{L}\) such that \([\mathbf{M}:\mathbf{L}] = n\) . (Observe firstly that there exists an irreducible polynomial \(f \in \mathbf{K}[\mathbf{K}]\) of degree \(n\) such that if \(\mathbf{N}\) is the splitting field of \(f\) , \(\operatorname{Gal}(\mathbf{N} / \mathbf{K})\) is isomorphic to \(\mathfrak{A}_n\) ; let \(y \in \Omega\) be a root of \(f\) in \(\Omega\) and let \(E = K(y)\) . Show that we necessarily have \(E \cap L = K\) and therefore \(L(E)\) is the desired field. Argue by contradiction, noting that the elements of \(E \cap L\) are A-constructible and that the relation \(E \cap L \neq K\) would imply the existence of a normal subgroup \(\Gamma\) of \(\mathfrak{A}_n\) such that \(1 < (\mathfrak{A}_n:\Gamma) < n\) , which would contradict the simplicity of \(\mathfrak{A}_n\) (I, p.~143, Ex. 16).
\end{enumerate}

\begin{enumerate}
\def\labelenumi{\arabic{enumi})}
\setcounter{enumi}{10}
\tightlist
\item
  \begin{enumerate}
  \def\labelenumii{\alph{enumii})}
  \tightlist
  \item
    Let E be a separable extension of a field K of finite degree and N the Galois extension of K generated by E. If \(\alpha\) is an element of N whose conjugates form a normal basis of N over K and if \(\beta = \operatorname{Tr}_{\mathbb{N}/\mathbb{E}}(\alpha)\) , show that \(\operatorname{E} = \operatorname{K}(\beta)\) .
  \end{enumerate}
\end{enumerate}

\begin{enumerate}
\def\labelenumi{\alph{enumi})}
\setcounter{enumi}{1}
\tightlist
\item
  Let N be a Galois extension of K of finite degree and let \(\alpha\) be an element of N whose conjugates form a normal basis over K. If L is a subextension of N which is Galois over K and if \(\beta = \mathrm{Tr}_{\mathbb{N}/\mathbb{L}}(\alpha)\) , show that the conjugates of \(\beta\) form a normal basis of L over K.
\item
  Let L, M be two Galois extensions of K of finite degree, such that \(L \cap M = K\) ; let \(\alpha\) (resp. \(\beta\) ) be an element of L (resp. M) whose conjugates over K form a normal basis of L (resp. M) over K; show that the conjugates over K of \(\alpha\beta\) form a normal basis of \(\mathrm{K}(\mathrm{L} \cup \mathrm{M})\) .
\end{enumerate}

\begin{enumerate}
\def\labelenumi{\arabic{enumi})}
\setcounter{enumi}{11}
\tightlist
\item
  Let K be an infinite field, N a Galois extension of degree n of K, \(f \in K[X]\) the minimal polynomial of an element \(\theta\) of N such that \(\mathbf{N} = \mathbf{K}(\theta)\) . The group \(\Gamma = \operatorname{Gal}(N/K)\) is allowed to operate on the polynomials in N{[}X{]} by operating on the coefficients of these polynomials.
\end{enumerate}

\begin{enumerate}
\def\labelenumi{\alph{enumi})}
\tightlist
\item
  Let \(g \in \mathbb{N}[\mathbf{X}]\) be the polynomial of degree \(\leqslant n - 1\) such that \(g(\theta) = 1\) , \(g(\sigma(\theta)) = 0\) for all \(\sigma \neq 1\) in \(\Gamma\) . For every subgroup \(\Delta\) of \(\Gamma\) let \(p_{\Delta} \in \mathbb{N}[\mathbf{X}]\) be the polynomial \(\det(\sigma\tau(g))_{\sigma \in \Delta, \tau \in \Delta}\) . Show that
\end{enumerate}

\[
p _ {\Delta} ^ {2} \equiv \sum_ {\sigma \in \Delta} \sigma (g) \mod f
\]

and deduce that \(p_{\Delta} \neq 0\) .

\begin{enumerate}
\def\labelenumi{\alph{enumi})}
\setcounter{enumi}{1}
\tightlist
\item
  Deduce that there exists an element \(\alpha \in K\) such that for every subgroup \(\Delta\) of \(\Gamma\) the conjugates of \(g(\alpha)\) over the subfield \(k(\Delta)\) form a normal basis of N over \(k(\Delta)\) .
\end{enumerate}

\begin{enumerate}
\def\labelenumi{\arabic{enumi})}
\setcounter{enumi}{12}
\item
  Let E be an algebraic extension of a field K, \(E_{s}\) the largest separable extension and \(E_{r}\) the largest p-radical extension of K contained in E. In an algebraic closure of K let N be the quasi-Galois extension of K generated by E; then the largest separable extension of K contained in N is the quasi-Galois extension \(N_{s}\) generated by \(E_{s}\) . For \(E_{r}\) to be the largest p-radical extension of K contained in N, it is necessary and sufficient that E should be separable over \(E_{r}\) (cf.~V, p.~152, Ex. 3).
\item
  Let \(\mathbf{E}\) be an algebraic extension of a field \(\mathbf{K}\) , \(\Gamma\) the group of \(\mathbf{K}\) -automorphisms of \(\mathbf{E}\) and \(\mathbf{S} \subset \mathbf{E}\) the field of invariants of \(\Gamma\) .
\end{enumerate}

\begin{enumerate}
\def\labelenumi{\alph{enumi})}
\item
  Show that for E to be quasi-Galois over K it is necessary and sufficient for S to be a p-radical extension of K.
\item
  Let \(S_{s}\) be the largest separable extension of K contained in S. Show that \(S_{s}\) is the smallest of the fields F such that \(K \subset F \subset E\) and E is a quasi-Galois extension of F.
\item
  Let \(E_{s}\) be the largest separable extension of K contained in E; show that \(\mathrm{E} = \mathrm{S}(\mathrm{E}_{s})\) (observe that no K-automorphism of E other than the identity leaves all the elements of \(\mathrm{S}(\mathrm{E}_{s})\) invariant).
\end{enumerate}

\begin{enumerate}
\def\labelenumi{\arabic{enumi})}
\setcounter{enumi}{14}
\item
  Let E be a Galois extension of finite degree n of a field K whose Galois group G is nilpotent. Let P be the set of prime numbers dividing n.~Then there exist Galois subextensions \(E_{p}\) of E ( \(p \in P\) ) such that \(\operatorname{Gal}(E_{p}/K)\) is a p-group for each \(p \in P\) , and such that the canonical homomorphism \(\otimes E_{p} \to E\) is an isomorphism.
\item
  Let \(a\) be an integer and let \(x_1, x_2, x_3, x_4\) be the roots of \(P(X) = X^4 - aX - 1\) in an algebraic closure of \(\mathbf{Q}\) . We identify the Galois group \(G\) of \(P\) with a group of permutations of \(\{x_1, \ldots, x_4\}\) . Suppose that \(Q(x_1)\) contains no quadratic subextension (cf.~Ex. 1 of \(V\) , p.~154).
\end{enumerate}

\begin{enumerate}
\def\labelenumi{\alph{enumi})}
\item
  Show that \(G = S_{4}\) or \(A_{4}\) . (Use Ex. 1 of V, p.~154: P is irreducible, hence G acts transitively; further \(\mathbf{Q}(x_{1})\) contains no quadratic subextensions, hence the stabilizer H of \(x_{1}\) is not contained in a subgroup of index 2 of G, whereas \([G:H] = 4\) ; deduce that G is not a 2-group, hence its order is divisible by 3.)
\item
  Let \(\delta = \prod_{1\leqslant i < j\leqslant 4}(x_i - x_j)\) and \(d = \delta^2\) . Show that \(d = 4^4 -3^3 a^4\) ; further \(s(\delta) = \varepsilon_s\delta\) for
\end{enumerate}

all \(s \in G\) . Deduce that \(G = \mathfrak{A}_4\) if \(\delta \in Q\) , and \(G = \mathfrak{S}_4\) otherwise. Show that \(\delta\) is not rational. (We have thus constructed an infinite family of equations of degree 4 over \(Q\) with Galois group \(\mathfrak{S}_4\) .)

* 17) Let \(f\) be an irreducible polynomial in \(\mathbf{Q}[\mathbf{X}]\) of prime degree \(p\) , having exactly two roots in \(\mathbf{C} - \mathbf{R}\) . Then the Galois group \(G\) of \(f\) is isomorphic to \(\mathfrak{S}_p\) . (G contains an element of order \(p\) and a transposition.) *

\begin{enumerate}
\def\labelenumi{\arabic{enumi})}
\setcounter{enumi}{17}
\tightlist
\item
  Let K be a field and F the field of rational fractions \(\mathrm{K}(\mathrm{T})\) . The group \(\mathrm{GL}_{2}(\mathrm{K})\) operates on F by the law \((\gamma f)(\mathrm{T}) = f\left(\frac{a\mathrm{T} + b}{c\mathrm{T} + d}\right)\) , for \(\gamma = \begin{pmatrix} a & b \\ c & d \end{pmatrix} \in \mathrm{GL}_{2}(\mathrm{K})\) and \(f \in F\) . The centre \(K^{*}\) . I of \(\mathrm{GL}_{2}(\mathrm{K})\) operates trivially, whence we obtain, by passage to quotients, a homomorphism of \(\mathrm{PGL}_{2}(\mathrm{K}) = \mathrm{GL}_{2}(\mathrm{K}) / \mathrm{K}^{*}\) . I into the automorphism group of the extension F of K. This homomorphism is bijective (V, p.~148, Ex. 11).
\end{enumerate}

\begin{enumerate}
\def\labelenumi{\alph{enumi})}
\tightlist
\item
  Let H be a subgroup of \(\mathrm{PGL}_{2}(\mathbf{K})\) of finite order h, and let E be the field of invariants of H, so that \(F = E(T)\) is a Galois extension of E with Galois group H (V, p.~66, Th. 3). Let
\end{enumerate}

\[
\mathrm{H} (\mathrm{X}) = \prod_ {s \in \mathrm{H}} (\mathrm{X} - s (\mathrm{T})) = \sum_ {i = 0} ^ {h} (- 1) ^ {i} \mathrm{S} _ {i} (\mathrm{T}) \mathrm{X} ^ {h - i}
\]

be the minimal polynomial of T over E. For each i such that \(0 \leqslant i \leqslant h\) , either \(S_{i}(T) \in K\) or \(E = K(S_{i}(T))\) . (Use Ex. 11, a) of V, p.~157.)

\begin{enumerate}
\def\labelenumi{\alph{enumi})}
\setcounter{enumi}{1}
\item
  For each \(s \in H\) let \(\begin{pmatrix} a_{s} & b_{s} \\ c_{s} & d_{s} \end{pmatrix}\) be a representative in \(GL_{2}(K)\) and let \(H_{0}\) be the subgroup of H consisting of those s for which \(c_{s} = 0\) ; let \(h_{0} = CardH_{0}\) and put \(R^{H}(T) = S_{h_{0}}(T)\) . Show that \(E = K(R^{H}(T))\) and that \(R^{H}(T) = P(T)/Q(T)^{h_{0}}\) where P is a polynomial of degree h and \(Q(T) = \prod_{s} (c_{s}T + d_{s})\) where s runs over a system of representatives of the cosets \(H_{0}s\) different from \(H_{0}\) ; the degree of Q is thus \((h/h_{0}) - 1\) .
\item
  Suppose that \(K^{*}\) contains a subgroup \(\mu_{n}\) of finite order n.~Let \(C_{n}\) be the image in \(PGL_{2}(K)\) of the subgroup of \(GL_{2}(K)\) consisting of the matrices \(\begin{pmatrix} a & 0 \\ 0 & 1 \end{pmatrix}, a \in \mu_{n}\) . Then \(C_{n}(X) = X^{n} - T^{n}\) and \(R^{C_{n}}(T) = (-1)^{n+1}T^{n}\) .
\item
  Let \(w\) be the element of \(\mathrm{PGL}_2(\mathbf{K})\) represented by \(\begin{pmatrix} 0 & 1 \\ 1 & 0 \end{pmatrix}\) . Then \(w^2 = 1\) and \(wsw^{-1} = s^{-1}\) if \(s\) is represented by a diagonal matrix. The subgroup \(\mathrm{D}_n\) of \(\mathrm{PGL}_2(\mathbf{K})\) generated by \(\mathrm{C}_n\) and \(w\) is therefore of order \(2n\) . We have \(\mathrm{D}_n(\mathbf{X}) = (\mathbf{X}^n - \mathrm{T}^n)\left(\mathbf{X}^n - \left(\frac{1}{\mathrm{T}}\right)^n\right) = \mathbf{X}^{2n} + (-1)\mathbf{R}^{\mathrm{D}_n(\mathrm{T})} + 1\) where \(\mathbf{R}^{\mathrm{D}_n(\mathrm{T})} = -((-T)^n + (-T)^{-n})\) .
\item
  Let S be the group of order 3 in \(\mathrm{PGL}_2(\mathbf{K})\) generated by the image of the matrix \(\left( \begin{array}{cc}0 & -1\\ 1 & 1 \end{array} \right)\) . We have \(S(X) = (X - T)\left(X - \frac{1}{1 - T}\right)\left(X - \frac{T - 1}{T}\right)\) and \(R^s (T) = \frac{T^3 - 3T + 1}{T^2 - 1}\) .
\item
  Suppose that K is a finite field with q elements. Let U be the image in \(\mathrm{PGL}_{2}(\mathbf{K})\) of the group of all matrices \(\begin{pmatrix}1 & 0 \\ b & 1\end{pmatrix}\) , \(b \in K\) . Then
\end{enumerate}

\[
\mathrm{U} (\mathrm{X}) = (\mathrm{X} - \mathrm{T}) ^ {q} - (\mathrm{X} - \mathrm{T}) = \mathrm{X} ^ {q} - \mathrm{X} - (\mathrm{T} ^ {q} - \mathrm{T}) \quad \text { and } \quad \mathrm{R} ^ {\mathrm{U}} (\mathrm{T}) = \mathrm{T} ^ {q} - \mathrm{T}.
\]

Let B be the group generated by \(C_{q-1}\) and U; it is of order \(q(q-1)\) . We have

\[
\mathrm{B} (\mathrm{X}) = \left(\mathrm{X} ^ {q} - \mathrm{X}\right) ^ {q - 1} - \left(\mathrm{T} ^ {q} - \mathrm{T}\right) ^ {q - 1} \quad \text { and } \quad \mathrm{R} ^ {\mathrm{H}} (\mathrm{T}) = \left(\mathrm{T} ^ {q} - \mathrm{T}\right) ^ {q - 1}.
\]

\begin{enumerate}
\def\labelenumi{\arabic{enumi})}
\setcounter{enumi}{18}
\item
  Let E be a Galois extension of a field K with Galois group G. For each \(x \in E\) denote by Gx the set of conjugates of x in E; this is a finite G-set. Let x, y be two elements of E. For the G-sets Gx and Gy to be isomorphic (that is, for a bijection \(Gx \rightarrow Gy\) to exist, compatible with the operation of G on the two sets), it is necessary and sufficient that the extensions \(\mathrm{K}(x)\) and \(\mathrm{K}(y)\) should be conjugate.
\item
  Let E be an algebraic extension of a field K such that every non-constant polynomial in K{[}X{]} has at least one root in E. Show that E is algebraically closed. (In an algebraic closure of E, hence of K, let F be an extension of finite degree of K. To show that \(F \subset E\) reduce to the case where F is quasi-Galois and then, by means of Prop. 13 of V, p.~76, to the case where F is either Galois or p-radical.)
\item
  Let \(f(\mathbf{X}) = \mathbf{X}^{n} + a_{1}\mathbf{X}^{n-1} + \cdots + a_{n}\) be a polynomial of Z{[}X{]} and p a prime number. Suppose that \(a_{i} \equiv 0 \pmod{p}\) , \(1 \leqslant i \leqslant n\) and that \(a_{n} \not\equiv 0 \pmod{p^{2}}\) . Show that \(f(\mathbf{X})\) is irreducible in Q{[}X{]}. (Let \(f(\mathbf{X}) = \mathrm{P}(\mathbf{X}) \mathrm{Q}(\mathbf{X})\) be a decomposition in Q{[}X{]}. Show that if P and Q are monic, they have integer coefficients. Reduce mod p.)
\end{enumerate}

* 22) Let \(m\) be an even integer \(>0\) , \(r\) an integer \(\geqslant 3\) and \(n_1 < \cdots < n_{r-2}\) , \(r-2\) even integers. Put \(g(\mathbf{X}) = (\mathbf{X}^2 + m)(\mathbf{X} - n_1) \ldots (\mathbf{X} - n_{r-2})\) .

\begin{enumerate}
\def\labelenumi{\alph{enumi})}
\item
  Put \(f(\mathbf{X}) = g(\mathbf{X}) - 2\) . Show that \(f(\mathbf{X})\) has at least \(r - 2\) real roots. Calculate the sum of the squares of the roots of \(f\) in \(\mathbf{C}\) . Deduce that for \(m \geqslant \sum n_i^2 / 2\) , \(f(\mathbf{X})\) is a polynomial with integer coefficients having \(r - 2\) real roots and two non-real conjugate roots.
\item
  Show that \(f(\mathbf{X})\) is irreducible in \(\mathbf{Q}[\mathbf{X}]\) . (Write \(f(\mathbf{X}) = \mathbf{X}' + \sum_{k=1}^{r} a_k \mathbf{X}'^{-k}\) ; show that
\end{enumerate}

\(a_{k}\) is even and that \(a_{r}\) is not divisible by 4; apply Ex. 21.)

\begin{enumerate}
\def\labelenumi{\alph{enumi})}
\setcounter{enumi}{2}
\item
  Show that when \(r\) is prime and \(m \geqslant \sum n_i^2 / 2\) , then the Galois group \(G\) of the splitting field of \(f\) is isomorphic to \(\mathfrak{S}_r\) (cf.~V, p.~158, Ex. 17).
\item
  Show that when \(r = 4\) and \(m \geqslant \sum n_i^2 / 2\) , the degree over \(\mathbf{Q}\) of the splitting field is equal to 8 or 24, and that it is equal to 8 if and only if \(n_1 = -n_2\) or also if and only if \(f(\mathbf{X})\) is of the form \(\mathrm{P}(\mathbf{X}^2)\) , where \(\mathbf{P}\) is a polynomial of degree 2. *
\end{enumerate}

\begin{enumerate}
\def\labelenumi{\arabic{enumi})}
\setcounter{enumi}{22}
\tightlist
\item
  Let \(k\) be a field of characteristic 2.
\end{enumerate}

\begin{enumerate}
\def\labelenumi{\alph{enumi})}
\tightlist
\item
  Let \(A_{1}, \ldots, A_{n}\) be indeterminates, K the field \(k(\mathbf{A}_{1}, \ldots, \mathbf{A}_{n})\) , E a splitting field of the polynomial \(P = X^{n} + A_{1}X^{n-1} + \cdots + A_{n} \in K[X]\) and \(X_{1}, \ldots, X_{n}\) the roots of P in E, so that \(E = k(X_{1}, \ldots, X_{n})\) (V, p.~21). We identify the Galois group of E/K with the symmetric group \(S_{n}\) acting by permutations of the variables \(X_{1}, \ldots, X_{n}\) (V, p.~59). Put
\end{enumerate}

\[
\beta \left(\mathrm{X} _ {1}, \dots , \mathrm{X} _ {n}\right) = \sum_ {i <   j} \mathrm{X} _ {i} / \left(\mathrm{X} _ {i} + \mathrm{X} _ {j}\right) \in \mathrm{E} \text { and } \mathrm{C} \left(\mathrm{A} _ {1}, \dots , \mathrm{A} _ {n}\right) = \sum_ {i <   j} \left(\mathrm{X} _ {i} \mathrm{X} _ {j}\right) / \left(\mathrm{X} _ {i} ^ {2} + \mathrm{X} _ {j} ^ {2}\right) \in \mathrm{K}.
\]

Then \(\beta \notin K\) and

\[
\boldsymbol {\beta} ^ {2} + \boldsymbol {\beta} + \mathbf {C} = 0.
\]

We have \(\sigma(\beta) = \beta\) for \(\sigma \in \mathfrak{A}_n\) and \(\sigma(\beta) = \beta + 1\) for \(\sigma \in \mathfrak{S}_n - \mathfrak{A}_n\) . Deduce that \(K(\beta)\) is the unique quadratic subextension of \(E/K\) .

\begin{enumerate}
\def\labelenumi{\alph{enumi})}
\setcounter{enumi}{1}
\tightlist
\item
  Let \(D \in Z[T_{1}, \ldots, T_{n}]\) be the discriminant of the polynomial \(X^{n} + T_{1}X^{n-1} + \cdots + T_{n}\) . Show that there exist Q, R, \(S \in Z[T_{1}, \ldots, T_{n}]\) with
\end{enumerate}

\[
\mathbf {D} = \mathbf {Q} ^ {2} - 4 \mathbf {R} + 8 \mathbf {S}
\]

and

\[
\mathrm{C} = \mathrm{C} (\mathrm{A} _ {1}, \dots , \mathrm{A} _ {n}) = \frac {\mathrm{R} (\mathrm{A} _ {1} , \dots , \mathrm{A} _ {n})}{\mathrm{Q} ^ {2} (\mathrm{A} _ {1} , \dots , \mathrm{A} _ {n})} = \frac {\mathrm{R} (\mathrm{A} _ {1} , \dots , \mathrm{A} _ {n})}{\mathrm{D} (\mathrm{A} _ {1} , \dots , \mathrm{A} _ {n})}.
\]

For every triple of elements U, V, W of \(\mathbf{Z}[T_1, \ldots, T_n]\) such that \(D = U^2 - 4V + 8W\) there exists \(M \in \mathbf{Z}[A_1, \ldots, A_n]\) such that

\[
C = \frac {V (A _ {1} , \dots , A _ {n})}{D (A _ {1} , \dots , A _ {n})} + \left(\frac {M}{D (A _ {1} , \dots , A _ {n})}\right) ^ {2} + \frac {M}{D (A _ {1} , \dots , A _ {n})}.
\]

\begin{enumerate}
\def\labelenumi{\alph{enumi})}
\setcounter{enumi}{2}
\tightlist
\item
  Let \(f = X^{n} + a_{1}X^{n-1} + \cdots + a_{n}\) be a separable polynomial of k{[}X{]} and G the Galois group of a splitting extension. Show that \((a_{1}, \ldots, a_{n})\) is substitutable in C and that the following conditions are equivalent
\end{enumerate}

\begin{enumerate}
\def\labelenumi{(\roman{enumi})}
\item
  every element of \(G\) induces an even permutation of the roots of \(f\) ;
\item
  there exists \(\alpha \in k\) such that \(C(a_{1},\ldots ,a_{n}) = \alpha^{2} + \alpha\) .
\end{enumerate}

* \(d\) ) Henceforth assume that \(k\) is finite and choose the elements U, V, W of \(\mathbf{Z}[T_1, \ldots, T_n]\) arbitrarily to satisfy the condition stated in \(b\) ).

Show that (i) and (ii) of \(c\) ) are equivalent to the following conditions:

\begin{enumerate}
\def\labelenumi{(\roman{enumi})}
\setcounter{enumi}{2}
\tightlist
\item
  \[
\operatorname{Tr} _ {k / \mathbf {F} _ {2}} (\mathbf {C}) = 0;
\]
\item
  \[
\mathrm{Tr} _ {k / \mathbf {F} _ {2}} \left(\frac {\mathbf {V} (a _ {1} , \dots , a _ {n})}{\mathbf {D} (a _ {1} , \dots , a _ {n})}\right) = 0.
\]
\end{enumerate}

\begin{enumerate}
\def\labelenumi{(\alph{enumi})}
\setcounter{enumi}{21}
\tightlist
\item
  The number of irreducible factors of \(f\) is congruent to \(n\) modulo 2 (use Ex. 22 of V, p.~165). *
\end{enumerate}

